\subsection{Deleting a corner fixed on one side}

\begin{lemma}[One-sided corner deletion]
\label{C:one-sided-corner-deletion}
Let $C$ be a finite category, let $a\ne b$, and suppose there are
arrows $u:b\to a$ and $v:a\to b$ such that
\[
 vu=e,\qquad e^2=e\ne1_b.
\]
Assume that $eg=g$ for every arrow $g:d\to b$ with $d\ne b$,
and that every unit of $\End_C(b)$ fixes every such incoming arrow.
Then the inclusion of the full subcategory $D=C\setminus\{b\}$
induces a homotopy equivalence $BD\simeq BC$.
The unit hypothesis holds in particular if the identity is the
only unit, or if $e$ is a zero of $\End_C(b)$.
There is a dual assertion with outgoing arrows in place of incoming
arrows.  No condition on outgoing arrows is needed in the stated
incoming version.
\end{lemma}
\begin{proof}
Put $A=\End_C(b)$ and write $A^\times$ for its unit group.
In a finite monoid, the nonunits form a two-sided ideal: a
one-sided inverse is an inverse, as follows by applying the
corresponding injective self-map of the finite monoid.
Define a functor $G:C\to D$ on objects by $G(b)=a$, fixing
every other object.  It fixes arrows of $D$ and is given on
the remaining arrows by
\begin{align*}
 G(h)&=1_a &&(h\in A^\times),\\
 G(h)&=uhv &&(h\in A\setminus A^\times),\\
 G(g)&=ug &&(g:d\to b,\ d\ne b),\\
 G(f)&=fv &&(f:b\to d,\ d\ne b).
\end{align*}
These formulas preserve identities.  We check every other
composition involving $b$.

For two nonunit endomorphisms $h,k$, their product is a nonunit
and
\[
 G(h)G(k)=uhvukv=uhekv=uhkv=G(hk),
\]
because $kv:a\to b$ is an incoming cross arrow fixed by $e$.
If one factor is a unit, its action on every incoming arrow is
the identity.  Thus $hkv=kv$ when $h$ is a unit, and $kv=v$
when $k$ is a unit.  These identities verify both mixed
endomorphism products; a product of two units is immediate.

Composing a nonunit $h$ after an incoming arrow $g$ gives
\[
 G(h)G(g)=uhvug=uheg=uhg=G(hg).
\]
The unit case uses $hg=g$.  Composing a nonunit $h$ before an
outgoing arrow $f$ gives
\[
 G(f)G(h)=fvuhv=fehv=fhv=G(fh),
\]
because $hv$ is incoming; the unit case uses $hv=v$.
For a composite through $b$ between retained objects,
\[
 G(f)G(g)=fvug=feg=fg.
\]
For a composite $gf:b\to b$ through a retained object,
$e(gf)=gf$.  It cannot be a unit, since multiplying that
identity by its inverse would imply $e=1_b$.  Consequently
\[
 G(g)G(f)=ugfv=G(gf).
\]
Compositions with other retained arrows follow directly by
associativity.  These cases exhaust all typed pairs, so $G$
is indeed a functor.

Let $J:D\to C$ be the inclusion.  Then $GJ=\id_D$.
The components $1_d$ for $d\ne b$ and $v$ at $b$ define a
natural transformation $JG\Rightarrow\id_C$.  For an incoming
$g$, its naturality equation is $vug=eg=g$; for an outgoing
$f$, it is the definition $G(f)=fv$.  For a nonunit $h$ it is
$vuhv=ehv=hv$, while for a unit it is $v=hv$.
Thus $BJ$ and $BG$ are homotopy inverses.
Finally, if $e$ is a zero of $A$ and $h$ is a unit, then
$hg=heg=eg=g$, which verifies the stated sufficient condition.
Opposition proves the dual assertion.
\end{proof}

\begin{corollary}[Arbitrary exits from the eleven-arrow strong pair]
\label{C:eleven-strong-arbitrary-exit}
Let $S$ be the eleven-arrow strong two-object category in
Proposition~\ref{C:eleven-strong-three-exit}.  Adjoin a sink
with trivial endomorphism monoid and any nonempty finite
compatible family of outgoing arrows from $S$.
Every rational degree-two square of the resulting classifying
space is zero.  In particular it cannot be a projective plane.
The dual assertion holds for an adjoined source.
\end{corollary}
\begin{proof}
If the two-dimensional row and column sets occur at different
objects, the balanced-tensor computation in the cited proposition
gives $H^2(BS;\mathbb Q)=0$.  Every degree-two class of the
whole act cone is then relative to $BS$ and pulls back from
the suspension of the action category.  Its square is zero.

Otherwise both two-dimensional sets occur at one object $a$,
with corner $P_5$, while the other object $b$ has corner
$\{1_b,e\}$.  Every incoming cross arrow at $b$ comes from
$a$ and belongs to the rectangular ideal, so is fixed by $e$.
The rectangular multiplication supplies $u,v$ with $vu=e$.
Lemma~\ref{C:one-sided-corner-deletion} deletes $b$ without
imposing any condition on the outgoing arrows to the new sink.
The resulting category is a finite act cone over $P_5$.
Theorem~\ref{C:two-row-act-cups} makes all its degree-two
squares zero, with no bound on the number of states.
\end{proof}
